% 08_desparasitar.tex
\begin{landscape}

	% =============================================================================
	% CONFIGURACIÓN
	% =============================================================================
	\setlength{\tabcolsep}{2pt}
	\setlength{\extrarowheight}{3pt}
	\fontsize{10pt}{10pt}\selectfont

	% =============================================================================
	% CÁLCULO DE ANCHOS DE COLUMNAS
	% =============================================================================
	\setlength{\availablewidth}{\linewidth}

	% Restamos (8 columnas * 2 lados) = 16 espacios para un ajuste matemático exacto
	\addtolength{\availablewidth}{-16\tabcolsep}

	% Restamos el ancho de las líneas verticales (9 líneas * 0.4pt estándar)
	\addtolength{\availablewidth}{-3.6pt}

	\newcommand{\tableheader}{\rowcolor{blue!10}
		\textbf{ESPE\-CIE} & \textbf{Nº DE CASOS} & \textbf{PROCEDIMIENTO} & \textbf{FÁRMACO} &
		\textbf{PRINCIPIO ACTIVO} & \textbf{DOSIS/VÍA DE AD\-MINISTRACIÓN} &
		\textbf{DURA\-CIÓN} & \textbf{OBSERVA\-CIO\-NES} \\
	}

	% =============================================================================
	% LONGTABLE CON CAPTION AL PRINCIPIO
	% =============================================================================
	\begin{longtable}{
		|C{0.07\availablewidth}          % ESPECIE
		|C{0.07\availablewidth}           % Nº CASOS
		|J{0.17\availablewidth}           % PROCEDIMIENTO
		|L{0.12\availablewidth}          % FÁRMACO
		|J{0.26\availablewidth}          % PRINCIPIO ACTIVO
		|L{0.16\availablewidth}          % DOSIS/VÍA
		|C{0.06\availablewidth}           % DURACIÓN
		|L{0.09\availablewidth}|         % OBSERVACIONES
	}

		% ──────────────── CAPTION COMO TÍTULO (solo en primera página) ────────────────
		\caption{\raggedright Registro de Desparasitaciones y aplicación de vitaminas realizadas durante el internado}
		\label{tab:desparasitaciones} \\
		\hline
		\tableheader
		\hline
		\endfirsthead

		% ──────────────── ENCABEZADO PARA PÁGINAS SIGUIENTES ────────────────
		\multicolumn{8}{l}{\normalsize\textbf{(Continuación)} Cuadro \ref{tab:desparasitaciones}. Registro de Desparasitaciones y aplicación de vitaminas realizadas durante el internado} \\[6pt]
		\hline
		\tableheader
		\hline
		\endhead

		% ──────────────── PIE DE PÁGINA (continuación) ────────────────
		\hline
		\multicolumn{8}{r}{\footnotesize\textit{continúa en la siguiente página \ldots}}
		\endfoot

		% ──────────────── PIE FINAL ────────────────
		%\hline
		\multicolumn{8}{l}{\footnotesize\textit{\textbf{Fuente}: Elaboración propia.}}
		\endlastfoot

		% =============================================================================
		% CONTENIDO DE LA TABLA (extraído del PDF '08_desparasitaciones.pdf')
		% =============================================================================

		% CANINO – Desparasitación externa
		\multirow{2}{*}{\textbf{Canino}} & \multirow{2}{*}{8} & \multirow{2}{*}{\parbox{\linewidth}{Desparasitación externa}} &
		Impacto & Cipermetrina, clorpirifos, citronelal & 1ml/1litro de agua & \multirow{2}{*}{1 mes} & \\
		\cline{4-6}
		& & & Galerfin & Clorfeniramina maleto & 0,2-2ml IV-IM & & \\
		\hline

		% CANINO – Desparasitación interna
		\multirow{3}{*}{\textbf{Canino}} & \multirow{3}{*}{67} & \multirow{3}{*}{\parbox{\linewidth}{Desparasitación interna}} &
		Levocan & Levamisol & 5 gotas/1kg Oral & \multirow{3}{*}{3 meses} & \\
		\cline{4-6}
		& & & Basken & Pamoato de pirantel, pamoato de oxantel & 1ml/1kg Oral & & \\
		\cline{4-6}
		& & & Vermic total & Praziquantel, pirantel pamoato, febantel & 1 tableta/10kg oral & & \\
		\hline

		% CANINO – Desparasitación interna-externa
		\multirow{2}{*}{\textbf{Canino}} & \multirow{2}{*}{58} & \multirow{2}{=}{Desparasitación interna y externa} &
		Ranger & Ivermectina 1\% & 1ml/30kg SC & \multirow{2}{*}{3 meses} & \\
		\cline{4-6}
		& & & Galerfin & Clorfeniramina maleto & 0,2-2ml IV-IM & & \\
		\hline

		% CANINO – Aplicación de vitaminas y minerales
		\multirow{2}{*}{\textbf{Canino}} & \multirow{2}{*}{125} & \multirow{2}{=}{Aplicación de vitaminas y minerales} &
		AminoLab Forte & Aminoácidos, dextrosa, cloruro de potasio, cloruro de sodio, cloruro de calcio, acetato de sodio, vitamina B1 B2 	B6 	B12 & 1ml/1kg IV & \multirow{2}{=}{\centering Cada 15 días} & \\
		\cline{4-6}
		& & & Complejo B & Vitamina B1, Vitamina B2, Vitamina B6, Vitamina B12, Vitamina B15 & 1ml/20kg IM-SC & & \\
		\hline

		% BOVINO – Desparasitación interna-externa
		\textbf{Bovino} & 160 & Desparasitación interna y externa & Ranger & Ivermectina 1\% & 1ml/50kg SC & 3-6 meses & \\
		\hline

		\rowcolor{green!10} % Verde claro para el total
		\textbf{TOTAL} & \textbf{418} \\
		\cline{1-2}

	\end{longtable}
	\vspace{-12pt}\normalsize
	% Resumen conciso de la tabla
	El cuadro registra las desparasitaciones y aplicaciones de vitaminas realizadas durante el internado. Se documentaron 418 casos: 8 desparasitaciones externas, 67 internas, 58 combinadas y 125 aplicaciones de vitaminas en caninos; 160 desparasitaciones combinadas en bovinos. Los tratamientos incluyen antiparasitarios como Ranger (ivermectina), Levocan (levamisol), Basken (pirantel), Vermic total (praziquantel), antihistamínicos como Galerfin, vitaminas como AminoLab Forte y Complejo B, administrados por vías oral, subcutánea, intramuscular o intravenosa durante 1 mes a 6 meses.

\end{landscape}
